\section{Related Work}

\label{sec:related}

\subsection{Decentralized Compute}

Golem~\cite{golem2016} pioneered decentralized compute with an orderbook model. Akash~\cite{akash2020} uses reverse auctions for cloud-like workloads. Render Network~\cite{rendernetwork2021} focuses on GPU rendering. Gensyn~\cite{gensyn2023} targets ML training verification. io.net~\cite{ionet2024} aggregates GPU clusters. ACI differs by providing a purpose-built blockchain with AI-native consensus and verification.

\subsection{Verifiable Computation}

TrueBit~\cite{teutsch2019truebit} introduced verification games for off-chain computation. zkEVM~\cite{zkevm2023} enables zero-knowledge proofs of EVM execution. RISC Zero~\cite{risczero2023} provides general-purpose ZK proofs. ACI's PoAI mechanism is specifically designed for AI workloads, combining TEE attestations with statistical verification rather than full replay or ZK proofs.

\subsection{Blockchain Consensus}

Narwhal-Tusk~\cite{danezis2022narwhal} introduced DAG-based consensus for high throughput. HotStuff~\cite{yin2019hotstuff} provided linear-communication BFT consensus. Sui~\cite{sui2023} combines DAG with object-based execution. ACI's dual-track consensus adapts these ideas for the specific requirements of compute marketplaces.

\subsection{AI Blockchain Integration}

Bittensor~\cite{bittensor2023} creates a decentralized network for AI model training. Ritual~\cite{ritual2024} brings AI inference on-chain. Modulus~\cite{modulus2024} provides ZK proofs for ML inference. ACI provides a more comprehensive infrastructure layer with purpose-built consensus, marketplace, and verification.

